%---------------------------------------------------------------------
\begin{abstract}
This manuscript reviews the scientific evidence and practitioner consensus relevant to rock climbing training across five domains: hangboarding, weekly training integration, supplementation, injury management, and recovery. Evidence is graded using an Evidence Strength Score (ESS 0--10) where ESS 10 = well-powered climbing RCT with functional outcomes and ESS 1--2 = anecdotal or mechanistic-only support; every actionable recommendation carries a status label (EVIDENCE\_BACKED, BIOPLAUSIBLE, FOLKLORE\_VERIFIED, or TO\_VERIFY).

For \textbf{hangboarding}, ten intervention trials (nine RCTs and one controlled study, combined $n \approx 230$), one retrospective cohort ($n = 527$), and one systematic review and meta-analysis \mbox{(Stien et~al.\ 2023, $n = 110$)} support three principal conclusions: (1) supplemental hangboarding 2--3$\times$/week consistently improves maximal finger strength (MFS) and grip endurance within 4--10 weeks in intermediate-to-advanced male climbers (meta-analytic SDM~= 0.41 for MFS, 95\%~CI 0.03--0.80; SDM~= 0.57 overall, 95\%~CI 0.24--0.91); (2) $\geq$80\% MFS protocols maximise MFS, 60--80\% MFS repeater protocols maximise stamina; (3) daily low-intensity hangs produced a positive but non-significant trend versus climbing-only controls ($+2.5\%$, $p = 0.12$). No trial has demonstrated a statistically significant climbing-grade improvement attributable to hangboarding alone; all effect sizes are on surrogate dynamometric measures.

For \textbf{weekly integration}, session sequencing rules are predominantly BIOPLAUSIBLE (ESS 2); the $\geq$48\,h inter-session minimum is a study-design convention, not an empirically validated prescription. For \textbf{supplementation}, no supplement has been tested with a functional climbing primary outcome; collagen + vitamin C and beta-alanine carry the strongest mechanistic and surrogate evidence for connective-tissue support and endurance tasks, respectively (ESS 3--7); caffeine showed null results in the only climbing-specific RCT. For \textbf{injury management}, A2 pulley conservative rehabilitation with pulley-protection splint (PPS) is the evidence-supported first-line approach (ESS 5; 88.4\% return to prior level at mean 8.8\,months in a prospective series); Grade I--III injuries are managed conservatively; surgery is reserved for Grade IVa/IVb or FLIP complications. The section additionally covers flexor tenosynovitis (second most common finger injury; favourable natural history, ESS 5), PIP capsulitis/synovitis (third most common; 72\% return with staged algorithm, ESS 5), lumbrical tears (climbing-specific; lumbrical stress test pathognomonic, ESS 5), adolescent physeal injuries (most common injury under age 16--17; pulley framework contraindicated), and NSAIDs (mechanistic caution for prolonged use; no climbing-specific trial). Experience-dependent risk is observationally established (ESS 5). For \textbf{recovery}, sleep $\geq$8\,h and strategic cold water immersion carry the strongest evidence base across sports, with no climbing-specific RCT for any modality.

The literature is uniformly limited by small samples (median PEDro 5--7/10), male-advanced cohorts, short follow-up ($\leq$10 weeks), and surrogate outcomes. Female-cohort data, long-duration trials with injury tracking, and any trial measuring grade improvement as a primary outcome are absent across all five domains.
\end{abstract}
